%% ============================================================
%% ICU PATIENT FACE MONITORING SYSTEM
%% UCS 503 -- SOFTWARE ENGINEERING PROJECT REPORT
%% ============================================================

\documentclass{tietreport}

%% ============================================================
%% PACKAGES
%% ============================================================

\usepackage{graphicx}
\usepackage{amsmath}
\usepackage{amssymb}
\usepackage{booktabs}
\usepackage{tabularx}
\usepackage{array}
\usepackage{float}
\usepackage{caption}
\usepackage{enumitem}
\usepackage{xcolor}
\usepackage{listings}
\usepackage{setspace}

\usepackage[
    nameinlink,
    noabbrev
]{cleveref}

\usepackage[
    backend=biber,
    style=alphabetic,
    sorting=nty,
    maxnames=5,
    minnames=3,
    maxalphanames=4,
    minalphanames=3,
    backref=true,
    doi=false,
    isbn=false,
    url=true,
    eprint=false
]{biblatex}

\addbibresource{references.bib}

%% ============================================================
%% GENERAL FORMATTING
%% ============================================================

\onehalfspacing

\setlist[itemize]{
    leftmargin=*,
    itemsep=3pt,
    topsep=4pt
}

\setlist[enumerate]{
    leftmargin=*,
    itemsep=4pt,
    topsep=4pt
}

\newcolumntype{Y}{
    >{\raggedright\arraybackslash}X
}

\newcolumntype{P}[1]{
    >{\raggedright\arraybackslash}p{#1}
}

\captionsetup{
    font=small,
    labelfont=bf,
    justification=centering
}

%% ============================================================
%% CODE LISTINGS
%% ============================================================

\lstdefinestyle{projectcode}{
    basicstyle=\ttfamily\small,
    breaklines=true,
    frame=single,
    columns=fullflexible,
    keepspaces=true,
    showstringspaces=false,
    numbers=left,
    numberstyle=\tiny,
    xleftmargin=2em,
    framexleftmargin=1.5em
}

\lstset{
    style=projectcode
}

%% ============================================================
%% PROJECT INFORMATION
%% ============================================================

\TitlePageHeader{
    UCS 503 -- Software Engineering Project Report
}

\ProjectTitle{
    ICU Patient Face Monitoring System
}

\ProjectSubTitle{
    A Computer-Vision-Based Non-Contact System for Continuous
    Facial and Ocular Monitoring of Intensive Care Unit Patients
}

\author{
    \texttt{1024240039} Harjot Singh\\
    \texttt{1024240103} Pahul Singh
}

\TitlePageSubText{
    BE Third Year, Data Science and Artificial Intelligence
}

\AdvisorName{
    Ms. Anushka
}

\TitlePageFooterText{
    \large Project Report\\[0.2cm]
    \bfseries Computer Science and Engineering Department\\[0.15cm]
    Thapar Institute of Engineering and Technology, Patiala
}

%% ============================================================
%% DOCUMENT
%% ============================================================

\begin{document}

%% ============================================================
%% TITLE PAGE
%% ============================================================

\maketitle

%% ============================================================
%% FRONT MATTER
%% ============================================================

\pagenumbering{roman}

%% ------------------------------------------------------------
%% ABSTRACT
%% ------------------------------------------------------------

\chapter*{Abstract}

\addcontentsline{toc}{chapter}{Abstract}

This report presents the design and implementation of the Face + Eye
component of an ICU Patient Monitoring System. The objective of the
component is to provide a non-contact computer-vision pipeline capable
of transforming camera observations of a patient into structured and
time-aware facial and ocular monitoring information.

The implemented pipeline begins with camera-frame acquisition and
progresses through face detection, facial landmark extraction, eye
landmark extraction, Eye Aspect Ratio (EAR) calculation, eye-state
classification, blink analysis, temporal processing, gaze estimation,
confidence estimation, global-value generation, alert generation, and
backend integration.

A key design principle is that individual video frames should not be
treated as independent monitoring observations. Eye measurements can be
affected by landmark noise, illumination, head movement, temporary
occlusion, and camera positioning. The system therefore uses temporal
history and duration-based reasoning to distinguish short eye closures
from persistent closure conditions.

The communication layer is implemented using a FastAPI backend. The
Face + Eye processing pipeline sends structured feature information
through an HTTP/JSON interface, while dashboard frames are transmitted
through a separate background sender. This separation reduces the effect
of dashboard streaming delays on the main processing pipeline.

The current project scope covers the Face + Eye processing pipeline and
its integration with the backend API. Full clinical validation,
patient-specific calibration, large-scale deployment, historical
storage, and production dashboard integration remain future work.

\noindent\textbf{Keywords:}
ICU monitoring, computer vision, facial landmarks, Eye Aspect Ratio,
blink detection, gaze estimation, temporal processing, FastAPI,
OpenCV, Python.

%% ------------------------------------------------------------
%% ABBREVIATIONS
%% ------------------------------------------------------------

\chapter*{List of Abbreviations}

\addcontentsline{toc}{chapter}{List of Abbreviations}

\begin{tabularx}{\textwidth}{@{}lY@{}}

\toprule
\textbf{Abbreviation} & \textbf{Meaning}\\
\midrule

API & Application Programming Interface\\
ASGI & Asynchronous Server Gateway Interface\\
EAR & Eye Aspect Ratio\\
FPS & Frames Per Second\\
HTTP & Hypertext Transfer Protocol\\
ICU & Intensive Care Unit\\
JSON & JavaScript Object Notation\\
REST & Representational State Transfer\\
TP & True Positive\\
FP & False Positive\\
FN & False Negative\\

\bottomrule

\end{tabularx}

%% ============================================================
%% TABLE OF CONTENTS
%% ============================================================

\tableofcontents

%% ============================================================
%% LIST OF FIGURES
%% ============================================================

\listoffigures

%% ============================================================
%% LIST OF TABLES
%% ============================================================

\listoftables

%% ============================================================
%% MAIN MATTER
%% ============================================================

\clearpage

\pagenumbering{arabic}

%% ============================================================
%% CHAPTER 1
%% ============================================================

\chapter{Introduction}

\section{Background and Context}

Intensive Care Units (ICUs) require continuous observation because the
condition of a patient can change rapidly. Conventional ICU monitoring
systems primarily rely on physiological sensors such as ECG
electrodes, pulse oximeters, blood-pressure monitors, and other medical
devices. These systems provide important physiological measurements
but do not directly capture visible facial and ocular behaviour.

Computer vision provides an opportunity to supplement conventional
monitoring through non-contact observation. A camera can continuously
capture the patient's face, while computer-vision algorithms can
identify facial landmarks and extract information from the eye region.
Vision-based patient monitoring has previously been studied as a
possible supplement to conventional monitoring systems
\cite{BDMS16,DMS19}.

The Face + Eye component developed as part of the ICU Patient
Monitoring System follows this approach. Rather than stopping at face
detection, the system progressively converts raw video frames into
structured monitoring information.

\begin{center}
\textbf{Camera}
$\rightarrow$
\textbf{Face Detection}
$\rightarrow$
\textbf{Landmarks}
$\rightarrow$
\textbf{Eye Features}
$\rightarrow$
\textbf{Temporal Processing}
$\rightarrow$
\textbf{Global Values}
$\rightarrow$
\textbf{Alerts}
$\rightarrow$
\textbf{Backend}
\end{center}

The system separates visual processing from communication. The
computer-vision layer focuses on extracting observations, while the
backend provides an interface through which the processed information
can be consumed by downstream components.

The system is intended as a supplementary monitoring and decision-support
system. It is not intended to replace professional medical judgement
or independently diagnose a patient's medical condition.

%% ------------------------------------------------------------

\section{Problem Statement}

Continuous visual observation of ICU patients can be challenging when
healthcare professionals must monitor several patients simultaneously.
Facial and ocular changes may occur continuously and may not always be
noticed immediately through manual observation.

The project therefore addresses the following problem:

\begin{quote}
\itshape
How can facial and ocular information obtained from a camera be
automatically transformed into structured, time-aware monitoring
information and made available through a backend service for
continuous ICU observation?
\end{quote}

The problem consists of several engineering tasks:

\begin{enumerate}
    \item detecting the patient's face;
    \item extracting facial landmarks;
    \item extracting eye-region landmarks;
    \item calculating eye-related geometric features;
    \item classifying observable eye states;
    \item analysing eye behaviour over time;
    \item estimating relative gaze behaviour;
    \item generating global monitoring values;
    \item generating observation-level alerts; and
    \item communicating processed information to a backend.
\end{enumerate}

%% ------------------------------------------------------------

\section{Motivation}

A single video frame is often insufficient for reliable behavioural
interpretation. A low EAR value can result from an actual blink,
incorrect landmarks, temporary occlusion, head movement, or poor image
quality.

Consequently, temporal processing is required to determine whether a
condition persists over multiple frames.

A second motivation is software modularity. The monitoring interface
should not need to know how facial landmarks were detected or how EAR
was calculated. A backend API establishes a clear abstraction boundary
between the computer-vision engine and downstream applications.

%% ------------------------------------------------------------

\section{Project Objectives}

The Face + Eye implementation is organised around the following
objectives:

\begin{enumerate}

    \item \textbf{Face Detection:}
    Identify the patient's face from incoming video frames.

    \item \textbf{Facial Landmark Detection:}
    Obtain structured facial landmark coordinates.

    \item \textbf{Eye Landmark Extraction:}
    Extract landmarks associated with the left and right eyes.

    \item \textbf{Eye Feature Extraction:}
    Calculate EAR and related geometric measurements.

    \item \textbf{Eye-State Analysis:}
    Distinguish observable states such as open, partially closed,
    and closed.

    \item \textbf{Temporal Analysis:}
    Analyse observations across multiple frames.

    \item \textbf{Global Value and Alert Generation:}
    Convert low-level measurements into compact monitoring values and
    observation-level alerts.

    \item \textbf{Backend Integration:}
    Expose processed information through a FastAPI-based
    communication layer.

\end{enumerate}

%% ------------------------------------------------------------

\section{Scope and Boundaries}

The current implementation covers the Face + Eye processing chain up
to the backend API boundary.

\begin{center}

\begin{tabular}{c}

Camera Input\\
$\downarrow$\\
Frame Acquisition\\
$\downarrow$\\
Face Detection\\
$\downarrow$\\
Facial Landmarks\\
$\downarrow$\\
Eye Features\\
$\downarrow$\\
Temporal Processing\\
$\downarrow$\\
Global Monitoring State\\
$\downarrow$\\
Alert Generation\\
$\downarrow$\\
FastAPI Backend

\end{tabular}

\end{center}

The system does not claim independent medical diagnosis or replacement
of professional clinical judgement.

%% ------------------------------------------------------------

\section{Project Contributions}

The main contributions of the implementation are:

\begin{itemize}
    \item modular decomposition of the computer-vision pipeline;
    \item facial and ocular feature extraction;
    \item EAR-based eye-state analysis;
    \item temporal reasoning for blink and closure behaviour;
    \item relative gaze estimation;
    \item confidence and global-state generation;
    \item observation-level alert generation;
    \item separation of processing and communication layers;
    \item HTTP/JSON backend integration; and
    \item asynchronous dashboard-frame transmission.
\end{itemize}

%% ============================================================
%% CHAPTER 2
%% ============================================================

\chapter{System Analysis and Requirements}

\section{System Context}

The Face + Eye component forms one processing component of the larger
ICU monitoring architecture. The camera provides visual input, the
computer-vision pipeline converts the input into structured
observations, and the backend provides these observations to downstream
applications.

The principal entities are:

\begin{itemize}
    \item \textbf{Camera:} Provides patient video frames.
    \item \textbf{Face + Eye Engine:} Performs computer-vision
    processing.
    \item \textbf{Backend:} Receives and exposes processed information.
    \item \textbf{Healthcare Staff:} Consume monitoring information
    and alerts.
    \item \textbf{Administrator:} Configures system parameters where
    applicable.
\end{itemize}

%% ------------------------------------------------------------

\section{Functional Requirements}

\begin{table}[H]
\centering
\caption{Functional requirements}
\label{tab:functional}

\begin{tabularx}{0.95\textwidth}{@{}P{2cm}Y@{}}

\toprule
\textbf{ID} & \textbf{Requirement}\\
\midrule

FR1 & The system shall acquire video frames from a camera or video
source.\\
FR2 & The system shall detect the patient's face.\\
FR3 & The system shall obtain facial landmarks.\\
FR4 & The system shall extract landmarks associated with both eyes.\\
FR5 & The system shall calculate eye-related geometric features.\\
FR6 & The system shall classify observable eye states.\\
FR7 & The system shall analyse eye behaviour temporally.\\
FR8 & The system shall detect prolonged eye-closure conditions.\\
FR9 & The system shall generate structured global monitoring values.\\
FR10 & The system shall transmit processed values to the backend.\\
FR11 & The system shall support dashboard-frame transmission.\\

\bottomrule

\end{tabularx}
\end{table}

%% ------------------------------------------------------------

\section{Non-Functional Requirements}

\begin{table}[H]
\centering
\caption{Non-functional requirements}
\label{tab:nonfunctional}

\begin{tabularx}{0.95\textwidth}{@{}P{3.2cm}Y@{}}

\toprule
\textbf{Property} & \textbf{Requirement}\\
\midrule

Performance &
The system should operate close to real-time.\\

Modularity &
Processing stages should remain independently maintainable.\\

Reliability &
Temporary communication failures should not terminate the main
processing pipeline.\\

Maintainability &
Components should have clearly separated responsibilities.\\

Scalability &
The architecture should permit future multi-patient deployment.\\

Security &
Patient information should be transmitted and stored using appropriate
controls.\\

Interpretability &
Outputs should represent observable computer-vision states rather
than unsupported clinical diagnoses.\\

\bottomrule

\end{tabularx}
\end{table}

%% ------------------------------------------------------------

\section{System Constraints}

The system can be affected by:

\begin{itemize}
    \item camera placement;
    \item illumination;
    \item patient head orientation;
    \item facial occlusion;
    \item landmark quality;
    \item camera-to-patient distance;
    \item computational resources; and
    \item backend/network availability.
\end{itemize}

%% ============================================================
%% CHAPTER 3
%% ============================================================

\chapter{Methodology and System Design}

\section{Overall Architecture}

The system follows a modular processing architecture. Each stage has
a specific responsibility and passes structured information to the next
stage.

\begin{center}

\begin{tabular}{c}

Camera\\
$\downarrow$\\
Frame Acquisition\\
$\downarrow$\\
Face Detection\\
$\downarrow$\\
Facial Landmark Detection\\
$\downarrow$\\
Eye Landmark Extraction\\
$\downarrow$\\
EAR and Eye Features\\
$\downarrow$\\
Eye-State and Blink Analysis\\
$\downarrow$\\
Temporal Processing\\
$\downarrow$\\
Gaze / Confidence Analysis\\
$\downarrow$\\
Global Monitoring State\\
$\downarrow$\\
Alert Generation\\
$\downarrow$\\
Backend Integration

\end{tabular}

\end{center}

%% ------------------------------------------------------------

\section{Context-Level Data Flow}

\begin{figure}[H]
    \centering
    \includegraphics[
        width=0.82\textwidth
    ]{figures/context-dfd.jpeg}

    \caption{
        Context-level Data Flow Diagram of the ICU Face Monitoring
        System.
    }

    \label{fig:context-dfd}
\end{figure}

The context-level DFD represents the monitoring system as a single
process interacting with the camera, healthcare staff, and
administrator.

%% ------------------------------------------------------------

\section{Level-1 Data Flow Diagram}

\begin{figure}[H]
    \centering
    \includegraphics[
        width=0.92\textwidth
    ]{figures/level1-dfd.jpeg}

    \caption{
        Level-1 Data Flow Diagram of the Face + Eye processing
        pipeline.
    }

    \label{fig:level1-dfd}
\end{figure}

The Level-1 DFD decomposes the system into face and landmark
processing, eye-feature extraction, temporal processing, global-value
generation, and alert generation.

%% ------------------------------------------------------------

\section{Use Case Design}

\begin{figure}[H]
    \centering
    \includegraphics[
        width=0.90\textwidth
    ]{figures/usecase-diagram.jpeg}

    \caption{
        Use Case Diagram of the ICU Face Monitoring System.
    }

    \label{fig:usecase}
\end{figure}

The use-case model identifies the Camera, Healthcare Staff, and Admin
as the principal actors.

%% ------------------------------------------------------------

\section{UML Class Design}

\begin{figure}[H]
    \centering
    \includegraphics[
        width=0.78\textwidth
    ]{figures/class-diagram.jpeg}

    \caption{
        UML Class Diagram of the Face + Eye processing pipeline.
    }

    \label{fig:class}
\end{figure}

The class-level design separates frame acquisition, face detection,
landmark extraction, eye processing, feature extraction, temporal
analysis, global-value generation, and alert handling.

%% ------------------------------------------------------------

\section{Activity Flow}

\begin{figure}[H]
    \centering
    \includegraphics[
        width=0.68\textwidth
    ]{figures/activity-diagram.jpeg}

    \caption{
        Activity diagram for per-frame processing and prolonged
        eye-closure alert generation.
    }

    \label{fig:activity}
\end{figure}

The activity flow shows that each frame is checked for a usable face
before eye features are calculated. Temporal information is updated
before alert generation.

%% ------------------------------------------------------------

\section{Sequence Flow}

\begin{figure}[H]
    \centering
    \includegraphics[
        width=0.72\textwidth
    ]{figures/sequence-diagram.jpeg}

    \caption{
        Sequence diagram showing the journey of a frame through the
        processing pipeline.
    }

    \label{fig:sequence}
\end{figure}

The sequence illustrates the movement of information from frame
acquisition through feature calculation, temporal processing, alert
evaluation, and monitoring delivery.

%% ------------------------------------------------------------

\section{Technical Stack}

\begin{table}[H]
\centering
\caption{Technology stack}
\label{tab:stack}

\begin{tabularx}{0.94\textwidth}{@{}P{3.3cm}Y@{}}

\toprule
\textbf{Technology} & \textbf{Role}\\
\midrule

Python & Primary implementation language.\\
OpenCV & Camera access and image processing.\\
MediaPipe Face Mesh & Facial landmark estimation.\\
NumPy & Numerical and geometric computation.\\
FastAPI & Backend API framework.\\
Uvicorn & ASGI application server.\\
Requests & HTTP communication from the Face + Eye pipeline.\\
Git & Version control and collaboration.\\

\bottomrule

\end{tabularx}
\end{table}

OpenCV is used for video acquisition and image processing, while
MediaPipe provides facial geometry for downstream eye analysis
\cite{KAGG19,Ope26,Goo26}.

%% ============================================================
%% CHAPTER 4
%% ============================================================

\chapter{Implementation}

\section{Project Structure}

The implementation is divided into separate functional modules.

\begin{lstlisting}[caption={Representative project structure}]

Face + Eye Module/
|
+-- main.py
+-- backend_sender.py
|
+-- camera/
|   +-- camera_stream.py
|
+-- landmarks/
|   +-- landmark_extractor.py
|   +-- landmark_groups.py
|
+-- eye/
|   +-- gaze_estimator.py
|   +-- gaze_tracker.py
|
+-- l2cs/
|   +-- model.py
|   +-- pipeline.py
|   +-- results.py
|
+-- temporal/
|   +-- feature_buffer.py
|
+-- episodes/
|   +-- episode_detector.py
|
+-- confidence/
|   +-- confidence_estimator.py
|
+-- clinical_metrics/
|   +-- indices.py
|
+-- visualization/
|   +-- overlay.py
|
+-- logging_system/
|   +-- distress_logger.py
|
+-- backend/
    +-- main.py

\end{lstlisting}

%% ------------------------------------------------------------

\section{Face Detection}

The first stage determines whether a usable face is present in the
incoming frame. Subsequent processing is performed only when a valid
face observation is available.

Face detection answers the question:

\begin{quote}
\itshape
Where is the patient's face in the current frame?
\end{quote}

Traditional real-time face detection approaches such as Viola--Jones
demonstrated the feasibility of efficient face detection
\cite{VJ04}.

If no usable face is detected, the frame is treated as a missing or
invalid observation rather than being forced through the eye-analysis
pipeline.

%% ------------------------------------------------------------

\section{Facial Landmark Detection}

After a face has been identified, facial landmarks are extracted. Each
landmark represents a geometric location on the face. Together, these
landmarks form a structured representation of facial geometry.

The principal purpose of this stage is to provide coordinates required
by downstream eye and facial analysis.

%% ------------------------------------------------------------

\section{Eye Landmark Extraction}

The left and right eye regions are extracted from the facial landmark
geometry. Both eyes are considered because one eye may occasionally
have reduced visibility because of head orientation, illumination, or
temporary occlusion.

The extracted eye landmarks provide the coordinates required for
geometric measurements.

%% ------------------------------------------------------------

\section{Eye Aspect Ratio}

The principal geometric feature used by the eye-analysis pipeline is
the Eye Aspect Ratio (EAR).

For six representative eye landmarks
$p_1,p_2,p_3,p_4,p_5,p_6$:

\begin{equation}
EAR =
\frac{
\lVert p_2-p_6\rVert
+
\lVert p_3-p_5\rVert
}{
2\lVert p_1-p_4\rVert
}
\label{eq:ear}
\end{equation}

where $\lVert\cdot\rVert$ represents Euclidean distance.

For an open eye, the vertical distances tend to be larger, producing a
higher EAR. As the eye closes, the vertical distances decrease and EAR
tends to become smaller.

EAR-based blink detection using facial landmarks is a recognised
computer-vision technique \cite{SC16}.

For the two eyes, the average EAR can be represented as:

\begin{equation}
EAR_{\mathrm{avg}}
=
\frac{
EAR_L+EAR_R
}{
2
}
\label{eq:average-ear}
\end{equation}

%% ------------------------------------------------------------

\section{Eye-State Classification}

The calculated EAR can be mapped to an observable eye state using a
configured threshold.

\begin{table}[H]
\centering
\caption{EAR-based eye-state interpretation}
\label{tab:eye-state}

\begin{tabularx}{0.75\textwidth}{@{}P{5cm}Y@{}}

\toprule
\textbf{EAR Condition} & \textbf{Eye State}\\
\midrule

EAR above configured threshold & Open\\
EAR near configured threshold & Partially closed\\
EAR below configured threshold & Closed\\

\bottomrule

\end{tabularx}
\end{table}

The threshold is an implementation parameter rather than a universal
medical constant.

%% ------------------------------------------------------------

\section{Blink Detection}

A blink is treated as a temporal transition:

\begin{center}

Open
$\rightarrow$
Closing
$\rightarrow$
Closed
$\rightarrow$
Opening
$\rightarrow$
Open

\end{center}

Rather than declaring a blink from a single frame, the system observes
eye measurements across time.

%% ------------------------------------------------------------

\section{Temporal Processing}

Facial and ocular measurements are treated as time-dependent
observations.

For a sequence:

\[
x_1,x_2,\ldots,x_n
\]

a moving average over the latest $N$ measurements is:

\begin{equation}
\bar{x}_t
=
\frac{1}{N}
\sum_{i=0}^{N-1}
x_{t-i}
\label{eq:moving-average}
\end{equation}

Temporal smoothing reduces the effect of isolated noisy observations.

%% ------------------------------------------------------------

\section{Gaze Estimation}

The system also incorporates gaze-related information.

Gaze is represented as relative behaviour rather than as an exact
clinical gaze measurement.

Possible states include:

\begin{itemize}
    \item left;
    \item right;
    \item upward;
    \item downward;
    \item centre; and
    \item uncertain.
\end{itemize}

%% ------------------------------------------------------------

\section{Confidence Estimation}

Confidence provides an indication of the reliability of the current
observation.

Relevant factors include:

\begin{itemize}
    \item face-detection quality;
    \item landmark quality;
    \item eye visibility;
    \item head pose;
    \item illumination; and
    \item temporary occlusion.
\end{itemize}

%% ------------------------------------------------------------

\section{Global Monitoring State}

The low-level outputs are consolidated into a structured monitoring
record.

A representative structure is:

\begin{lstlisting}
{
    "eye_state": "open",
    "left_ear": 0.31,
    "right_ear": 0.30,
    "average_ear": 0.305,
    "blink_count": 2,
    "closure_duration": 0.0,
    "gaze_state": "center",
    "confidence": 0.94,
    "alert_state": "normal"
}
\end{lstlisting}

%% ------------------------------------------------------------

\section{Alert Generation}

The alert layer converts temporally processed information into
observation-level monitoring events.

For prolonged eye closure:

\begin{equation}
Alert =
\begin{cases}
1, & \text{if closure duration} > T_c,\\
0, & \text{otherwise}.
\end{cases}
\label{eq:alert}
\end{equation}

where $T_c$ represents the configured closure-duration threshold.

\begin{table}[H]
\centering
\caption{Observation-level alert categories}
\label{tab:alerts}

\begin{tabularx}{0.92\textwidth}{@{}P{4cm}Y@{}}

\toprule
\textbf{Alert} & \textbf{Description}\\
\midrule

Prolonged Eye Closure &
Eye closure persists beyond the configured duration.\\

Repeated Closure &
Multiple closure events occur in a recent period.\\

Unusual Blink Pattern &
Blink behaviour differs from configured criteria.\\

Low Confidence &
Face/eye observations are considered unreliable.\\

Face Not Detected &
A usable face cannot currently be identified.\\

\bottomrule

\end{tabularx}
\end{table}

%% ============================================================
%% CHAPTER 5
%% ============================================================

\chapter{Backend Integration}

\section{Integration Objective}

The backend provides a communication boundary between the Face + Eye
processing engine and downstream monitoring applications.

The architecture separates:

\begin{enumerate}
    \item \textbf{Processing Layer:} Video handling, face detection,
    landmarks, eye features, temporal processing, gaze, confidence,
    global values, and alerts.

    \item \textbf{Communication Layer:} HTTP requests, JSON payloads,
    frame uploads, endpoint handling, and response generation.
\end{enumerate}

%% ------------------------------------------------------------

\section{FastAPI Backend}

The backend is implemented using FastAPI and is served using Uvicorn.

The local development server is started using:

\begin{lstlisting}[language=bash]
uvicorn backend.main:app --host 0.0.0.0 --port 8000 --reload
\end{lstlisting}

The service is accessible locally through:

\begin{lstlisting}
http://127.0.0.1:8000
\end{lstlisting}

The interactive API documentation is available through:

\begin{lstlisting}
http://127.0.0.1:8000/docs
\end{lstlisting}

%% ------------------------------------------------------------

\section{Feature Communication}

The Face + Eye processing pipeline sends structured monitoring
information to:

\begin{lstlisting}
POST /api/face-eye/features
\end{lstlisting}

The request body contains the processed feature information in JSON
format.

%% ------------------------------------------------------------

\section{Dashboard Frame Communication}

Dashboard frames are transmitted separately through:

\begin{lstlisting}
POST /api/stream/face
\end{lstlisting}

Frames are encoded as JPEG before transmission.

%% ------------------------------------------------------------

\section{Asynchronous Frame Delivery}

Dashboard frame transmission is separated from the primary processing
path.

The sender maintains only the latest frame. A background thread waits
for new frames and uploads them to the backend.

This prevents old frames from accumulating in an unbounded queue and
reduces the effect of streaming delays on the main processing loop.

%% ------------------------------------------------------------

\section{Backend API}

The backend receives Face + Eye feature information and dashboard
frames.

A simplified implementation is:

\begin{lstlisting}[language=Python]
from fastapi import FastAPI, UploadFile, File
from fastapi.responses import StreamingResponse
import io

app = FastAPI()

face_eye_data = {}
posture_data = {}
latest_frame = None


@app.post("/api/face-eye/features")
def receive_face_eye(data: dict):

    global face_eye_data

    face_eye_data = data

    return {"received": True}


@app.get("/api/face-eye/features")
def get_face_eye():

    return face_eye_data


@app.post("/posture")
def receive_posture(data: dict):

    global posture_data

    posture_data = data

    return {"received": True}


@app.get("/posture")
def get_posture():

    return posture_data


@app.get("/fusion")
def fusion():

    return {
        "face_eye": face_eye_data,
        "posture": posture_data
    }


@app.post("/api/stream/face")
async def receive_frame(
    frame: UploadFile = File(...)
):

    global latest_frame

    latest_frame = await frame.read()

    return {"received": True}


@app.get("/api/stream/face")
def get_frame():

    if latest_frame is None:
        return {"error": "No frame available"}

    return StreamingResponse(
        io.BytesIO(latest_frame),
        media_type="image/jpeg"
    )
\end{lstlisting}

%% ------------------------------------------------------------

\section{End-to-End Integration Flow}

The completed communication architecture is:

\begin{center}

\begin{tabular}{c}

Camera\\
$\downarrow$\\
Face + Eye Processing\\
$\downarrow$\\
Global Monitoring Record\\
$\downarrow$\\
Feature JSON + JPEG Frame\\
$\downarrow$\\
FastAPI Backend\\
$\downarrow$\\
Monitoring Application

\end{tabular}

\end{center}

%% ------------------------------------------------------------

\section{Failure Handling}

Backend failures should not terminate the primary computer-vision
pipeline.

The feature sender catches HTTP request exceptions, while the frame
sender operates in a separate daemon thread.

This provides separation between computer-vision processing, clinical
feature communication, and dashboard frame streaming.

%% ============================================================
%% CHAPTER 6
%% ============================================================

\chapter{Testing and Evaluation}

\section{Testing Methodology}

Testing is organised at multiple levels.

\begin{table}[H]
\centering
\caption{Testing methodology}
\label{tab:testing}

\begin{tabularx}{0.95\textwidth}{@{}P{3.4cm}Y@{}}

\toprule
\textbf{Test Area} & \textbf{Representative Tests}\\
\midrule

Face Detection &
Normal illumination, low illumination, different distances,
face orientation, temporary occlusion.\\

Eye Landmarks &
Open eyes, closed eyes, blinking, partial visibility, head movement.\\

EAR Features &
Open versus closed eyes and left/right eye consistency.\\

Temporal Processing &
Short blink versus sustained closure and noisy frames.\\

Alert Generation &
Correct triggering, alert clearing, and avoidance of excessive
alerts.\\

Backend &
Server startup, endpoint accessibility, JSON validation, invalid
requests, and communication failures.\\

Streaming &
JPEG encoding, frame upload, latest-frame replacement, and streaming
failure isolation.\\

End-to-End &
Camera frame through processing, global state, alert generation,
and backend delivery.\\

\bottomrule

\end{tabularx}
\end{table}

%% ------------------------------------------------------------

\section{Performance Metrics}

The principal performance metrics are FPS, precision, recall, F1,
latency, and detection accuracy.

\begin{equation}
FPS =
\frac{
\text{Number of processed frames}
}{
\text{Processing time}
}
\end{equation}

\begin{equation}
Precision =
\frac{TP}{TP+FP}
\end{equation}

\begin{equation}
Recall =
\frac{TP}{TP+FN}
\end{equation}

\begin{equation}
F1 =
2
\frac{
Precision \times Recall
}{
Precision + Recall
}
\end{equation}

Latency represents the time between frame acquisition and the
availability of the corresponding processed information.

%% ------------------------------------------------------------

\section{Evaluation Targets}

\begin{table}[H]
\centering
\caption{Proposed evaluation targets}
\label{tab:targets}

\begin{tabularx}{0.80\textwidth}{@{}P{5cm}Y@{}}

\toprule
\textbf{Metric} & \textbf{Target}\\
\midrule

Processing latency &
Below approximately 50 ms per processing step.\\

Face detection accuracy &
Above 95\% on selected test data.\\

Blink detection accuracy &
Above 90\% on labelled sequences.\\

Throughput &
Approximately 15 FPS or higher.\\

\bottomrule

\end{tabularx}
\end{table}

These values are evaluation targets and must not be presented as
measured results unless experimentally validated.

%% ------------------------------------------------------------

\section{Backend Integration Validation}

The backend integration can be validated through the following
sequence:

\begin{enumerate}
    \item Start the FastAPI application using Uvicorn.
    \item Open \texttt{http://127.0.0.1:8000/docs}.
    \item Confirm that the Face + Eye endpoints are registered.
    \item Start the Face + Eye pipeline.
    \item Confirm that feature POST requests reach the backend.
    \item Confirm that frame POST requests reach the streaming endpoint.
    \item Query the fusion endpoint.
    \item Verify that Face + Eye information is present in the returned
    data.
\end{enumerate}

%% ------------------------------------------------------------

\section{Current Results}

The implemented work establishes the processing and communication
pipeline from camera input through the backend boundary.

The transformation achieved is:

\begin{center}

Camera Frame
$\rightarrow$
Face Information
$\rightarrow$
Landmarks
$\rightarrow$
Eye Measurements
$\rightarrow$
Temporal Information
$\rightarrow$
Global Monitoring State
$\rightarrow$
Alert Information
$\rightarrow$
Backend Data

\end{center}

A complete numerical accuracy evaluation requires labelled test data,
defined ground truth, repeatable conditions, and systematic measurement.

%% ============================================================
%% CHAPTER 7
%% ============================================================

\chapter{Discussion}

\section{Engineering Significance}

The principal engineering contribution is the transformation of raw
visual input into a structured monitoring representation.

The system creates an abstraction boundary around the processed
monitoring record, improving maintainability, testability, integration,
modularity, and future extensibility.

%% ------------------------------------------------------------

\section{Importance of Temporal Reasoning}

Temporal reasoning is a central design decision.

A single low EAR measurement may result from a normal blink, landmark
noise, head movement, poor illumination, temporary occlusion, or
sustained eye closure.

Using recent history and duration allows the system to distinguish
isolated measurements from persistent conditions.

%% ------------------------------------------------------------

\section{Processing and Communication Separation}

The processing engine and backend are deliberately separated.

The processing layer is responsible for video acquisition, face
detection, landmark extraction, eye features, temporal analysis, gaze,
confidence, global values, and alerts.

The backend is responsible for receiving structured data, serving
processed data, receiving frames, and exposing information to
downstream clients.

%% ------------------------------------------------------------

\section{Asynchronous Frame Streaming}

Dashboard frame transmission is deliberately separated from the
clinical feature path.

Only the newest frame is retained for transmission. This prevents a
slow network connection from producing an unbounded backlog of old
frames.

%% ------------------------------------------------------------

\section{Limitations}

The current system has several limitations:

\begin{itemize}
    \item sensitivity to poor illumination;
    \item sensitivity to large head-pose changes;
    \item possible occlusion from hands or medical equipment;
    \item dependency on camera placement;
    \item dependency on landmark quality;
    \item fixed thresholds rather than personalised baselines;
    \item relative rather than clinically calibrated gaze estimation;
    \item absence of large-scale clinical validation; and
    \item dependency on backend/network availability.
\end{itemize}

The output should therefore be considered supplementary monitoring
information rather than an independent medical diagnosis.

%% ============================================================
%% CHAPTER 8
%% ============================================================

\chapter{Conclusion and Future Work}

\section{Summary}

The Face + Eye component of the ICU Patient Monitoring System
implements a modular computer-vision pipeline for non-contact
observation of facial and ocular behaviour.

The processing chain begins with camera input and progresses through
face detection, facial landmarks, eye landmarks, EAR-based feature
extraction, eye-state analysis, blink analysis, temporal processing,
gaze estimation, confidence estimation, global-value generation,
alert generation, and backend integration.

%% ------------------------------------------------------------

\section{Key Findings}

\begin{enumerate}

    \item \textbf{Temporal processing improves robustness.}

    \item \textbf{Modularity improves maintainability.}

    \item \textbf{Structured API communication simplifies integration.}

    \item \textbf{Asynchronous streaming protects the main pipeline.}

    \item \textbf{The system is supplementary and requires validation
    before clinical deployment.}

\end{enumerate}

%% ------------------------------------------------------------

\section{Future Work}

\begin{enumerate}

    \item Patient-specific calibration.

    \item Improved temporal modelling.

    \item Improved gaze estimation.

    \item Multi-camera deployment.

    \item Historical monitoring.

    \item Scalable backend architecture.

    \item Advanced alert prioritisation.

    \item Systematic evaluation using labelled sequences.

    \item Production dashboard integration.

    \item Security and privacy controls.

\end{enumerate}

%% ------------------------------------------------------------

\section{Final Remarks}

The project demonstrates that a modular computer-vision system can
transform camera observations into structured, time-aware monitoring
information and expose that information through a backend service.

The strength of the implementation lies in the integration of multiple
processing stages rather than any single feature.

The system should be considered a supplementary monitoring and
decision-support component. Appropriate technical, clinical, privacy,
and security validation is required before deployment in a real ICU
environment.

%% ============================================================
%% BIBLIOGRAPHY
%% ============================================================

\printbibliography[
    heading=bibintoc,
    title={Bibliography}
]

%% ============================================================
%% END
%% ============================================================

\end{document}